\originalsection{第7次习题课}
\sourceinfo{MIT 18.04, Spring 2018, Recitation 7; 题面 \nolinkurl{mit18_04s18_recit7-handout.pdf}, PDF第1页; 解答 \nolinkurl{mit18_04s18_recit7-solutions.pdf}, PDF第1页. 课程作者 Jeremy Orloff, 页内署名 Vishesh Jain, 2018, CC BY-NC-SA 4.0}
\begin{exercise}
\textbf{原题 Rec7.1。}\label{ass:rec7-1}
求 $f(z)=(z+1)/(z^3(z^2+1))$ 在环域 $0<|z|<1$ 中以0为中心的 Laurent 级数.
\end{exercise}
\begin{solution}
编者补充 (AI): 官方解答仅指向讲义例7.22. 因 $|z|<1$, 几何级数 $1/(1+z^2)=\sum_{k=0}^\infty(-1)^kz^{2k}$ 收敛, 故
\[
f(z)=\sum_{k=0}^\infty(-1)^k(z^{2k-3}+z^{2k-2})
=z^{-3}+z^{-2}-z^{-1}-1+z+z^2-z^3-z^4+\cdots.
\]
乘以 $z^{-3}(z+1)$ 不改变正半径内的收敛性, 所以该式在所给环域成立, 并在每个紧子环域一致收敛. 主部为 $z^{-3}+z^{-2}-z^{-1}$.
\end{solution}
\begin{exercise}
\textbf{原题 Rec7.2。}\label{ass:rec7-2}
求 $f(z)=1/(z(z-1))$ 在以下区域以0为中心的 Laurent 级数: 2.1. $0<|z|<1$; 2.2. $1<|z|<\infty$.
\end{exercise}
\begin{solution}
编者补充 (AI): 官方解答仅指向讲义例7.23. 在第一个环域, 以 $z$ 作几何级数变量得 $f(z)=-z^{-1}/(1-z)=-\sum_{k=0}^\infty z^{k-1}=-z^{-1}-1-z-z^2-\cdots$. 在第二个环域, 改以 $1/z$ 为变量, 得 $f(z)=z^{-2}/(1-z^{-1})=\sum_{k=0}^\infty z^{-k-2}=z^{-2}+z^{-3}+\cdots$. 两个展开所在环域不同, 不要求 Laurent 系数相同; 中间的奇点 $z=1$ 阻止把第一个级数延伸到第二个环域.
\end{solution}
\begin{exercise}
\textbf{原题 Rec7.3。}\label{ass:rec7-3}
设 $f$ 在单位圆盘内解析, 对所有 $w\in\{x+\ii0:-0.2\le x\le0.2\}$ 有 $f(w)=5$. 求 $f(\ii/2)$. 若改为对所有 $w\in\{x-3\ii/4:-0.2\le x\le0.2\}$ 有 $f(w)=5$, 答案改变吗?
\end{exercise}
\begin{solution}
据官方解答翻译补全; 第二条线段的情况原解答遗漏, 此处为编者补充 (AI). $g=f-5$ 在连通单位圆盘内解析. 第一条线段上的零点以0为域内聚点, 恒等定理给出 $g\equiv0$, 故 $f(\ii/2)=5$. 第二条线段全部位于圆盘内, 因 $|x-3\ii/4|\le\sqrt{0.2^2+(3/4)^2}<1$. 零点以 $-3\ii/4$ 为域内聚点, 同一定理仍使 $f\equiv5$, 答案不变.
\end{solution}
\begin{exercise}
\textbf{原题 Rec7.4。}\label{ass:rec7-4}
求微分方程 $f'(x)=f(x)+2$, $f(0)=0$ 的幂级数解.
\end{exercise}
\begin{solution}
编者补充 (AI): 官方解答仅指向讲义例7.24. 设 $f(x)=\sum_{n\ge0}a_nx^n$, 初值给 $a_0=0$. 比较常数项得 $a_1=a_0+2=2$, 比较 $x^n$ 项 ($n\ge1$) 得 $(n+1)a_{n+1}=a_n$. 归纳得到 $a_n=2/n!$ ($n\ge1$), 所以 $f(x)=2\sum_{n=1}^\infty x^n/n!=2(\ee^x-1)$. 该级数收敛半径无穷大, 可逐项求导, 所得函数满足方程与初值. 递推唯一确定所有系数, 因而它也是唯一的解析解.
\end{solution}
